%; whizzy paragraph -pdf xpdf -latex ./whizzypdfptex.sh
%; whizzy-paragraph "^\\\\begin{frame}"
% latex beamer presentation.
% platex, latex-beamer でコンパイルすることを想定。 

%     Tokyo Debian Meeting resources
%     Copyright (C) 2009 Junichi Uekawa
%     Copyright (C) 2009 Nobuhiro Iwamatsu

%     This program is free software; you can redistribute it and/or modify
%     it under the terms of the GNU General Public License as published by
%     the Free Software Foundation; either version 2 of the License, or
%     (at your option) any later version.

%     This program is distributed in the hope that it will be useful,
%     but WITHOUT ANY WARRANTY; without even the implied warreanty of
%     MERCHANTABILITY or FITNESS FOR A PARTICULAR PURPOSE.  See the
%     GNU General Public License for more details.

%     You should have received a copy of the GNU General Public License
%     along with this program; if not, write to the Free Software
%     Foundation, Inc., 51 Franklin St, Fifth Floor, Boston, MA  02110-1301 USA

\documentclass[cjk,dvipdfmx,12pt]{beamer}
\usetheme{Tokyo}
\usepackage{monthlypresentation}

%  preview (shell-command (concat "evince " (replace-regexp-in-string "tex$" "pdf"(buffer-file-name)) "&")) 
%  presentation (shell-command (concat "xpdf -fullscreen " (replace-regexp-in-string "tex$" "pdf"(buffer-file-name)) "&"))
%  presentation (shell-command (concat "evince " (replace-regexp-in-string "tex$" "pdf"(buffer-file-name)) "&"))

%http://www.naney.org/diki/dk/hyperref.html
%日本語EUC系環境の時
\AtBeginDvi{\special{pdf:tounicode EUC-UCS2}}
%シフトJIS系環境の時
%\AtBeginDvi{\special{pdf:tounicode 90ms-RKSJ-UCS2}}

\title{僕の疑問を解決してください}
%\subtitle{第75回 2011年4月度}
\author{まえだこうへい mkouhei@debian.or.jp\\IRC nick: mkouhei}
\date{2011年4月16日}
\logo{\includegraphics[width=8cm]{image200607/openlogo-light.eps}}

\begin{document}

\frame{\titlepage{}}

\section{}

\begin{frame}
 \frametitle{最近、パッケージメンテナンスを再チャレンジし始めた}
  \begin{enumerate}
   \item 最近興味をもったいくつかのPythonのソフトウェア
   \item apt-cache searchで探したけど無い
   \item PyPiに登録されている
   \item 自分でDebianパッケージにしてみた
   \item 公式パッケージにしてしまおう
   \item スポンサーを岩松さんにお願いした←イマココ
  \end{enumerate}
 {\large べ、別に会長になったから、じゃないんだからね！}
\end{frame}

\begin{frame}{そこでいくつかの疑問}
 \begin{itemize}
  \item WNPPに入ってないけどBTSってどうすんの？
  \item debian/copyrightについての疑問
  \item ドキュメントが無い場合の疑問
  \item 複数に分けるべき？
 \end{itemize}
\end{frame}

\begin{frame}{WNPPに入って無いけどBTSってどうすんの？}
\begin{itemize}
 \item 今回パッケージをつくったのは、
       \begin{itemize}
	\item *diagシリーズ\footnote{blockdiag,
       sphinxcontlib-blockdiag, seqdiag, actdiag, netdiag}
	\item blockdiagが依存するfuncparserlib
	\item MessagePackのPythonバインディング\footnote{msgpack-python}
       \end{itemize}
 \item WNPPには入ってない
 \item こういう場合はBTSに登録するの？
 \item mentors.debian.netに登録すればよいだけ?
       \end{itemize}
\end{frame}

\begin{frame}{debian/copyrightについての疑問}
 \begin{itemize}
  \item 次のライセンスは/usr/share/common-licensesにあるのでパスを書けば
	良いという認識
	\begin{itemize}
	 \item msgpack-pythonはApache Software License
	 \item sphinxcontrib-blockdiagはBSD　License
	\end{itemize}
  \item funcparserlibはMITライセンスなので、dh\_makeでx11を指定すればよ
	いだけ？(common-licensesディレクトリにはない)
  \item 他のdiagシリーズは、Python Software Foundation License(PSFL)。Debianの
	mainには含まれている
	\footnote{http://www.debian.org/legal/licenses/}けど、この場合は
	PSFL\footnote{http://www.python.org/download/releases/2.5.2/license/}
	の全文をdebian/copyrightにコピペすればいいの？
	\end{itemize}
\end{frame}

\begin{frame}{ドキュメントがない場合の疑問}
 \begin{itemize}
  \item こういう場合は、自分で書くの？アップストリームに書いてもらうの？
  \item 自分で書く場合は、アップストリームにフィードバックして取り込んで
	もらうもん？
  \item アップストリームへの依頼は直接やればいいの？どこか経由で依頼した
	方がよいの？
 \end{itemize}
\end{frame}

\begin{frame}{複数パッケージに分けるべき？}
 \begin{itemize}
  \item blockdiag以外の*diagシリーズは、sphinxcontrib-*diagも含まれてい
	る。
  \item こういう場合は、パッケージとしては別にするべき？
 \end{itemize}
\end{frame}

\begin{frame}{お礼}
\begin{itemize}
 \item 疑問に答えてくれてありがとー
 \item 頑張ってやっていきますー
\end{itemize}
\end{frame}
\end{document}
